\documentclass[11pt]{article}

% ================================================================
% Packages
% ================================================================
\usepackage{amsmath,amssymb,amsthm}
\usepackage{mathrsfs}
\usepackage[shortlabels]{enumitem}
\usepackage{booktabs}
\usepackage{array}
\usepackage{microtype}
\usepackage{xcolor}
\usepackage[colorlinks=true,linkcolor=blue!60!black,
 citecolor=green!40!black,urlcolor=blue!60!black]{hyperref}

% ================================================================
% Theorem environments
% ================================================================
\newtheorem{theorem}{Theorem}[section]
\newtheorem{proposition}[theorem]{Proposition}
\newtheorem{lemma}[theorem]{Lemma}
\newtheorem{corollary}[theorem]{Corollary}
\theoremstyle{definition}
\newtheorem{definition}[theorem]{Definition}
\newtheorem{construction}[theorem]{Construction}
\newtheorem{computation}[theorem]{Computation}
\newtheorem{example}[theorem]{Example}
\theoremstyle{remark}
\newtheorem{remark}[theorem]{Remark}
\newtheorem{conjecture}[theorem]{Conjecture}

% ================================================================
% Macros
% ================================================================
\newcommand{\cA}{\mathcal{A}}
\newcommand{\cC}{\mathcal{C}}
\newcommand{\cG}{\mathcal{G}}
\newcommand{\cH}{\mathcal{H}}
\newcommand{\cM}{\mathcal{M}}
\newcommand{\Walg}{\mathcal{W}}
\newcommand{\barB}{\bar{B}}
\newcommand{\fg}{\mathfrak{g}}
\newcommand{\fsl}{\mathfrak{sl}}
\newcommand{\Vir}{\mathrm{Vir}}
\newcommand{\Heis}{\mathrm{Heis}}
\newcommand{\FP}{\mathrm{FP}}
\newcommand{\diag}{\mathrm{diag}}
\newcommand{\cross}{\mathrm{cross}}
\newcommand{\scal}{\mathrm{scal}}
\newcommand{\grav}{\mathrm{grav}}

\DeclareMathOperator{\Aut}{Aut}
\DeclareMathOperator{\av}{av}
\DeclareMathOperator{\ord}{ord}

\numberwithin{equation}{section}

% ================================================================
\begin{document}

\title{Multi-weight genus expansion and cross-channel corrections}

\author{Raeez Lorgat}
\date{}

\maketitle

\begin{abstract}
For a modular Koszul chiral algebra whose strong generators
have a single conformal weight, the genus-$g$ free energy
satisfies the scalar formula
$F_g = \kappa \cdot \lambda_g^{\FP}$ at all genera
(UNIFORM-WEIGHT).
This paper proves that the scalar formula fails for
multi-weight algebras at genus $g \geq 2$ and computes the
failure explicitly.
The full genus expansion decomposes as
$F_g = \kappa \cdot \lambda_g^{\FP}
+ \delta F_g^{\cross}$ (ALL-WEIGHT),
where $\delta F_g^{\cross}$ is a graph sum over
mixed-channel boundary strata of
$\overline{\cM}_{g,0}$.
The cross-channel correction vanishes identically at
genus~$1$, on the uniform-weight locus, and for all
free-field algebras.
For $\Walg_3$ we compute
$\delta F_2 = (c{+}204)/(16c)$ over the $7$ stable graphs of
$\overline{\cM}_{2,0}$ and
$\delta F_3$ over all $42$ stable graphs of
$\overline{\cM}_{3,0}$; the cross-channel correction is
already comparable to the scalar part at genus~$3$.
A universal closed-form $N$-formula for
$\delta F_2(\Walg_N)$ is proved, with the Virasoro algebra
($N=2$) as the unique point of vanishing.
The scalar and cross-channel sectors have
distinct resurgent structures: the scalar instanton action is
the universal $A_{\scal} = (2\pi)^2$, while the cross-channel
growth exceeds it at large central charge.
The cross-channel dominance is the quantitative expression
of $E_1$ primacy: the ordered bar complex
$B^{\ord}(\cA) = T^c(s^{-1}\bar{\cA})$ resolves all
channels simultaneously, and the cross-channel correction
measures the information lost by the coinvariant projection
$\av \colon \fg^{E_1} \to \fg^{\mathrm{mod}}$.
\end{abstract}

\tableofcontents


% ================================================================
% Section 1: Introduction
% ================================================================
\section{Introduction}\label{sec:introduction}

\subsection{The scalar formula and its failure}

The modular Koszul programme associates to each chiral algebra
$\cA$ a modular characteristic $\kappa(\cA)$ and a
Maurer--Cartan element $\Theta_{\cA}$ in the convolution dg~Lie
algebra $\fg^{\mathrm{mod}}_{\cA}$.
The genus-$g$ free energy extracts the weight-$g$ component of
the MC element through the stable-graph expansion on
$\overline{\cM}_{g,n}$.
When $\cA$ has strong generators of a single conformal weight
(the \emph{uniform-weight} case: Heisenberg, Virasoro, affine
Kac--Moody at any single current weight), the genus expansion
collapses to the scalar formula
\begin{equation}\label{eq:intro-scalar}
 F_g(\cA)
 \;=\;
 \kappa(\cA) \cdot \lambda_g^{\FP}
 \qquad\text{for all } g \geq 1
 \quad\text{(UNIFORM-WEIGHT)},
\end{equation}
% AP1: kappa from census. Heis: k. Vir: c/2. KM: dim(g)(k+h^v)/(2h^v). W_N: c*(H_N - 1).
where
$\lambda_g^{\FP}
 = (2^{2g-1}{-}1)|B_{2g}|/(2^{2g-1}\cdot(2g)!)$
is the Faber--Pandharipande intersection number
$\int_{\overline{\cM}_g} \lambda_g$
($\lambda_1^{\FP} = 1/24$,
$\lambda_2^{\FP} = 7/5760$,
$\lambda_3^{\FP} = 31/967680$).

The scalar formula encodes a profound universality: the genus-$g$
free energy is a single number $\kappa(\cA)$ multiplied by a
universal topological constant $\lambda_g^{\FP}$.
All information about the OPE, the representation category, and
the symmetry algebra is compressed into the single invariant
$\kappa$.

The question forced by the structure is whether this universality
survives the passage to multiple generators at distinct conformal
weights.
The $\Walg_3$ algebra has two generators:
$T$ at conformal weight $2$ and $W$ at conformal weight $3$.
Its modular characteristic is
% AP1: kappa(W_N) = c*(H_N - 1). H_3 = 1 + 1/2 + 1/3 = 11/6. H_3 - 1 = 5/6.
% kappa(W_3) = c * 5/6. Checks: N=2 -> H_2 - 1 = 1/2, kappa(W_2) = c/2 = kappa_Vir. Good.
$\kappa(\Walg_3) = 5c/6$,
and the scalar formula predicts
$F_2(\Walg_3) = (5c/6) \cdot (7/5760) = 7c/6912$.
The actual genus-$2$ free energy is
\begin{equation}\label{eq:intro-w3-genus2}
 F_2(\Walg_3)
 \;=\;
 \underbrace{\frac{7c}{6912}}_{\kappa \cdot \lambda_2^{\FP}
 \;\text{(UNIFORM-WEIGHT)}}
 \;+\;
 \underbrace{\frac{c+204}{16c}}_{\delta F_2^{\cross}
 \;\text{(ALL-WEIGHT)}}\,.
\end{equation}
The scalar formula fails. The correction
$\delta F_2^{\cross} = (c{+}204)/(16c)$ is nonzero for
all finite central charge $c > 0$ and resolves the
multi-generator universality problem in the negative.

\subsection{Scope and main results}

This paper develops the theory of the cross-channel correction
$\delta F_g^{\cross}$ from first principles and computes it
in the simplest non-trivial cases. The main results are:

\begin{enumerate}[label=(\arabic*)]
\item The decomposition
 $F_g = \kappa \cdot \lambda_g^{\FP}
 + \delta F_g^{\cross}$ (ALL-WEIGHT)
 holds for every modular Koszul chiral algebra at every genus
 $g \geq 1$; the diagonal part is universal and the
 cross-channel part is the sole obstruction to the scalar
 formula (Theorem~\ref{thm:multi-weight-decomposition}).

\item The cross-channel correction vanishes identically at
 genus~$1$ (single-edge argument), on the uniform-weight
 locus (channel symmetry), and for all free-field algebras
 (triple redundancy: shadow-tower collapse, off-diagonal
 metric elimination, ghost-number conservation).

\item For $\Walg_3$ at genus $2$:
 $\delta F_2 = (c{+}204)/(16c)$, computed graph by graph
 over all $7$ stable graphs of
 $\overline{\cM}_{2,0}$, of which $3$ contribute nonzero
 cross-channel terms (Computation~\ref{comp:w3-genus2}).

\item For $\Walg_3$ at genus $3$:
 $\delta F_3 = (5c^3 + 3792c^2
 + 1{,}149{,}120\,c + 217{,}071{,}360)/(138{,}240\,c^2)$,
 computed over all $42$ stable graphs, of which $34$
 contribute (Computation~\ref{comp:w3-genus3}).
 The cross-channel correction is comparable to the scalar
 part: the ratio
 $\delta F_3/(\kappa \cdot \lambda_3^{\FP})
 \to 42/31 \approx 1.35$ as $c \to \infty$.

\item A universal closed-form $N$-formula for the genus-$2$
 gravitational cross-channel correction of $\Walg_N$
 (Proposition~\ref{prop:universal-N-formula}), with
 $N = 2$ (Virasoro) as the unique vanishing point.

\item The cross-channel instanton hierarchy: the scalar
 sector has universal instanton action
 $A_{\scal} = (2\pi)^2$ (Gevrey-$1$), while the
 cross-channel growth at large $c$ exceeds the scalar
 growth (Proposition~\ref{prop:instanton-hierarchy}).

\item The $E_1$ interpretation: the cross-channel correction
 measures the information in the ordered bar complex
 $B^{\ord}(\cA)$ that is killed by the coinvariant
 projection to the symmetric bar $B^{\Sigma}(\cA)$
 (\S\ref{sec:e1-interpretation}).
\end{enumerate}

\subsection{Conventions}
% FOUR OBJECTS:
% 1. B(A) = bar coalgebra = T^c(s^{-1} A-bar) with deconcatenation + twist
% 2. A^i = H^*(B(A)) = dual coalgebra (Koszul cohomology of bar)
% 3. A^! = ((A^i)^v) = dual ALGEBRA (linear or Verdier dual)
% 4. Z^der_ch(A) = derived chiral center = Hochschild cochains = bulk

The bar complex is $B(\cA) = T^c(s^{-1}\bar{\cA})$ with
$\bar{\cA} = \ker(\varepsilon)$ the augmentation ideal,
$|s^{-1}v| = |v| - 1$ (desuspension lowers degree), and
deconcatenation coproduct.
All grading is cohomological ($|d| = +1$).
The ordered bar $B^{\ord}(\cA)$ is the primitive $E_1$-chiral
coalgebra; the symmetric bar $B^{\Sigma}(\cA)$ is its
$\Sigma_n$-coinvariant shadow.
The modular characteristic $\kappa(\cA)$ is
$\av(r(z))$ at degree $2$, where $r(z)$ is the classical
$r$-matrix and $\av \colon \fg^{E_1} \to \fg^{\mathrm{mod}}$
is the (lossy) coinvariant projection.
The Koszul complementarity sum is $K(\cA) = \kappa(\cA) + \kappa(\cA^!)$:
$K = 0$ for Kac--Moody and free fields, $K = 13$ for Virasoro,
$K = 250/3$ for $\Walg_3$.

Every displayed formula involving $F_g$ or
$\mathrm{obs}_g$ carries one of the two tags:
\textbf{(UNIFORM-WEIGHT)} when the formula holds under the
hypothesis that all strong generators share a single conformal
weight, or \textbf{(ALL-WEIGHT)} when the formula holds
unconditionally (possibly with a cross-channel correction term).


% ================================================================
% Section 2: The cross-channel correction
% ================================================================
\section{The cross-channel correction}\label{sec:cross-channel}

\subsection{Multi-channel graph sum}

\begin{construction}[Multi-channel graph sum]\label{constr:multi-channel}
Let $\cA$ be a modular Koszul chiral algebra of rank $r$
with strong generators $\varphi_1, \ldots, \varphi_r$ at
conformal weights $h_1, \ldots, h_r$ and diagonal
Zamolodchikov metric
$\eta = \diag(\eta_1, \ldots, \eta_r)$.
Fix a genus $g \geq 2$.
For each stable graph $\Gamma \in \cG_{g,0}$ with edge
set $E(\Gamma)$ and vertex set $V(\Gamma)$:

\emph{Channel assignment.}
A function $\sigma \colon E(\Gamma) \to \{1, \ldots, r\}$.
There are $r^{|E(\Gamma)|}$ such assignments.

\emph{Multi-channel amplitude.}
\begin{equation}\label{eq:multi-channel-amplitude}
 A_{\Gamma}(\sigma, \cA)
 \;:=\;
 \prod_{e \in E(\Gamma)} \eta^{\sigma(e)\sigma(e)}
 \;\cdot\;
 \prod_{v \in V(\Gamma)}
 V_{g(v), n(v)}\!\bigl(\sigma|_{H(v)}\bigr),
\end{equation}
where $\eta^{ii} = 1/\eta_i$ is the inverse Zamolodchikov
metric in channel $i$ (the propagator; each edge is diagonal
in the channel basis; weight~$1$ always),
$V_{g(v), n(v)}(\sigma|_{H(v)})$ is the genus-$g(v)$
amplitude at vertex $v$ with $n(v)$ half-edges, and
$g(v)$ is the genus label of $v$.

\emph{Full graph sum.}
\begin{equation}\label{eq:full-graph-sum}
 F_g(\cA)
 \;=\;
 \sum_{\Gamma \in \cG_{g,0}}
 \frac{1}{|\Aut(\Gamma)|}
 \sum_{\sigma \colon E(\Gamma) \to \{1,\ldots,r\}}
 A_{\Gamma}(\sigma, \cA).
\end{equation}

\emph{Diagonal/cross decomposition.}
An assignment $\sigma$ is \emph{diagonal} if
$\sigma(e) = i$ for all $e$ (constant on all edges);
otherwise it is \emph{mixed}.
The diagonal sum gives
$F_g^{\diag}(\cA) = \kappa(\cA) \cdot \lambda_g^{\FP}$
(UNIFORM-WEIGHT: per-channel universality; each single-channel
component has uniform weight).
The mixed sum gives $\delta F_g^{\cross}(\cA)$.
\end{construction}

\subsection{The decomposition theorem}

\begin{theorem}[Multi-weight genus expansion]\label{thm:multi-weight-decomposition}
Let $\cA$ be a modular Koszul chiral algebra with
strong generators $\varphi_1, \ldots, \varphi_r$ of conformal
weights $h_1, \ldots, h_r$, per-channel modular characteristics
$\kappa_1, \ldots, \kappa_r$,
and total modular characteristic
$\kappa(\cA) = \sum_{i=1}^r \kappa_i$.
\begin{enumerate}[label=\textup{(\roman*)}]
\item \emph{Per-channel universality.}
 The diagonal contribution to the genus-$g$ free energy satisfies
 \begin{equation}\label{eq:diagonal-universality}
 F_g^{\diag}(\cA)
 \;=\;
 \kappa(\cA) \cdot \lambda_g^{\FP}
 \qquad\text{for all } g \geq 1
 \quad\text{\textup{(}UNIFORM-WEIGHT scalar lane\textup{)}}.
 \end{equation}

\item \emph{Cross-channel decomposition.}
 The full genus-$g$ free energy decomposes as
 \begin{equation}\label{eq:all-weight-decomposition}
 F_g(\cA)
 \;=\;
 \kappa(\cA) \cdot \lambda_g^{\FP}
 \;+\;
 \delta F_g^{\cross}(\cA)
 \qquad\text{for all } g \geq 1
 \quad\text{\textup{(}ALL-WEIGHT\textup{)}},
 \end{equation}
 where $\delta F_g^{\cross}(\cA)$ is the graph sum over
 mixed-channel boundary graphs of $\overline{\cM}_{g,0}$
 \textup{(}Construction~\textup{\ref{constr:multi-channel}}\textup{)}.

\item \emph{Genus-$1$ universality.}
 $\delta F_1^{\cross} = 0$ for all $\cA$:
 the unique genus-$1$ boundary graph \textup{(}figure-eight\textup{)}
 has a single edge, admitting only pure channel assignments.

\item \emph{Uniform-weight universality.}
 If $h_1 = h_2 = \cdots = h_r$, then
 $\delta F_g^{\cross} = 0$ for all $g \geq 1$.

\item \emph{$R$-matrix independence.}
 The cross-channel correction
 $\delta F_g^{\cross}$ is independent of the
 Givental $R$-matrix: genus-$0$ boundary vertices receive
 no $R$-contribution, and genus-$1$ vertices have vanishing
 $R$-correction for degree reasons.

\item \emph{Free-field exactness.}
 If $\cA$ is a free-field algebra \textup{(}purely quadratic
 genus-$0$ OPE: $m_k = 0$ for $k \geq 3$\textup{)}, then
 $\delta F_g^{\cross}(\cA) = 0$ for all $g \geq 1$
 \textup{(}ALL-WEIGHT free-field exact case\textup{)}.
\end{enumerate}
\end{theorem}

\begin{proof}
(i) Each edge $e$ of a stable graph $\Gamma$ carries a single
channel assignment $\sigma(e) \in \{1, \ldots, r\}$.
The diagonal graph sum restricts to constant assignments
$\sigma(e) = i$ for all $e$.
On such an assignment, the amplitude factors as
$A_{\Gamma}(\sigma_i, \cA) = A_{\Gamma}(\cA_i)$,
the single-channel amplitude for the rank-$1$ algebra $\cA_i$
with curvature $\kappa_i$.
Summing over all $i$ and all graphs:
\[
 F_g^{\diag}(\cA)
 \;=\;
 \sum_{i=1}^r F_g(\cA_i)
 \;=\;
 \sum_{i=1}^r \kappa_i \cdot \lambda_g^{\FP}
 \;=\;
 \kappa(\cA) \cdot \lambda_g^{\FP}
 \quad\text{(UNIFORM-WEIGHT)},
\]
where the second equality uses single-channel universality:
each single-channel component has a single generator and
hence uniform weight.

\smallskip\noindent
(ii) The full graph sum runs over all channel assignments
$\sigma \colon E(\Gamma) \to \{1, \ldots, r\}$.
The diagonal assignments give $F_g^{\diag}$; the mixed
assignments give $\delta F_g^{\cross}$.

\smallskip\noindent
(iii) At genus $1$, the unique boundary graph is the
figure-eight (one vertex of genus $0$, one self-loop,
$|E| = 1$).
A single edge has only $r$ channel assignments, each pure.
No mixed-channel assignments exist on a single-edge graph.

\smallskip\noindent
(iv) When all conformal weights coincide, the channel
permutation symmetry forces every mixed-channel amplitude
to equal the corresponding diagonal amplitude; the
difference vanishes.

\smallskip\noindent
(v) The genus-$0$ vertices receive transferred operations
from the cyclic minimal model, which are independent of
the Givental $R$-matrix by construction.
Genus-$1$ vertices on $\overline{\cM}_{1,1}$ have vanishing
$R$-correction because the $R$-matrix acts in degree $\geq 2$
on the state space and the genus-$1$ vertex amplitude requires
only degree-$0$ and degree-$1$ data.

\smallskip\noindent
(vi) For free fields, three independent mechanisms each
force $\delta F_g^{\cross} = 0$:

\emph{Shadow-tower collapse.}
Free fields have $m_k = 0$ for $k \geq 3$, so every
genus-$0$ vertex is binary.
No higher-contact vertex couples generators of distinct
conformal weights at a common vertex.
Only diagonal assignments survive.

\emph{Off-diagonal metric elimination.}
The OPE pairing $g^{ab}$ is block-diagonal with respect to
conformal weight: $\varphi_a(z)\varphi_b(w) \sim 0$ for
$h_a \neq h_b$ forces $g^{ab} = 0$.
Each edge carries a propagator
$P^{ab} = g^{ab} \cdot d\log E(z,w)$ (weight~$1$),
so $P^{ab} = 0$ for cross-weight edges.
Every mixed-channel graph has at least one such edge and
therefore vanishes.

\emph{Ghost-number conservation.}
Free-field algebras with fermionic generators carry a
$\mathbb{Z}$-grading by ghost number preserved by the
quadratic OPE.
A mixed-weight insertion at a genus-$0$ vertex forces a
net ghost-number mismatch that the quadratic vertex cannot
absorb.

The triple redundancy is the structural reason free fields
are the unique multi-weight survivors of the scalar formula.
\end{proof}

\begin{remark}[The propagator ratio]\label{rem:propagator-ratio}
For a rank-$2$ algebra with curvatures $\kappa_1, \kappa_2$
and Zamolodchikov metric entries $\eta_1, \eta_2$, the
\emph{propagator ratio}
\begin{equation}\label{eq:propagator-ratio}
 \varrho_{12}
 \;:=\;
 \frac{\eta^{22}}{\eta^{11}}
 \;=\;
 \frac{\eta_1}{\eta_2}
\end{equation}
controls the strength of the cross-channel correction.
When $\varrho_{12} = 1$, the system is effectively
single-channel and $\delta F_g^{\cross} = 0$.
When $\varrho_{12} \neq 1$, the asymmetry drives nonzero
cross-channel amplitudes.
For $\Walg_3$: $\eta^{TT} = 2/c$, $\eta^{WW} = 3/c$,
so $\varrho_{TW} = 3/2 \neq 1$.
\end{remark}


% ================================================================
% Section 3: W_3 at genus 2
% ================================================================
\section{$\Walg_3$ at genus $2$: the seven stable graphs}
\label{sec:w3-genus2}

\subsection{The stable graph catalogue}

The moduli space $\overline{\cM}_{2,0}$ has seven
stable graph types.
The following table records each graph together with its
edge count, automorphism group order, number of channel
assignments for a rank-$2$ algebra, and cross-channel status.

\begin{center}
\small
\renewcommand{\arraystretch}{1.35}
\begin{tabular}{llcccc}
\toprule
Label & Name & $|E|$ & $|\Aut|$
 & Assign.\ & Cross status \\
\midrule
A & smooth & $0$ & $1$ & $1$ & none (no edges) \\
B & figure-eight & $1$ & $2$ & $2$ & pure only \\
C & separating & $1$ & $2$ & $2$ & genus-$1$ universal \\
D & sunset & $2$ & $8$ & $4$ & \textbf{mixed $\neq 0$} \\
E & bridge-loop & $2$ & $2$ & $4$ & \textbf{mixed $\neq 0$} \\
F & theta & $3$ & $12$ & $8$ & \textbf{mixed $\neq 0$} \\
G & dumbbell & $3$ & $8$ & $8$ & genus-$1$ universal \\
\bottomrule
\end{tabular}
\end{center}

Graphs A, B, C, G contribute zero cross-channel correction:
\begin{itemize}
\item Graph~A (smooth irreducible genus-$2$) has no edges.
\item Graph~B (figure-eight: one vertex of genus $0$, one
 self-loop) has a single edge; only pure assignments exist.
\item Graph~C (separating: two genus-$1$ vertices joined by
 a bridge) has one edge.
 Each genus-$1$ factor is universal by
 Theorem~\ref{thm:multi-weight-decomposition}(iii),
 so the clutching product is a product of
 genus-$1$-universal quantities.
\item Graph~G (dumbbell: two self-loops connected by a bridge,
 each self-loop a single-edge genus-$1$ component) has
 the same genus-$1$ universal structure.
\end{itemize}

Graphs D, E, F carry the entire cross-channel correction.
Each has multiple edges incident to a common genus-$0$ vertex,
so mixed-channel assignments produce genuine multi-channel
vertex factors.

\subsection{The computation}

\begin{computation}[$\Walg_3$ genus-$2$ cross-channel]\label{comp:w3-genus2}
The $\Walg_3$ algebra has rank $2$ with channels
$T$ (weight $2$, $\eta_{TT} = c/2$) and
$W$ (weight $3$, $\eta_{WW} = c/3$).
The inverse Zamolodchikov metrics are
$\eta^{TT} = 2/c$ and $\eta^{WW} = 3/c$;
the per-channel modular characteristics are
% AP1: kappa(W_3) = 5c/6. kappa_T = c/2 (Virasoro channel). kappa_W = c/3 (W-channel).
$\kappa_T = c/2$ and $\kappa_W = c/3$, with total
$\kappa(\Walg_3) = 5c/6$.
The $WW$ OPE contains the composite quasi-primary
$\Lambda = {:}TT{:} - \tfrac{3}{10}\partial^2 T$
with structure constant $16/(22{+}5c)$.

The three contributing graphs give:

\emph{Sunset \textup{(}graph D\textup{)}.}
Two self-loops at a single genus-$0$ vertex with $n = 4$.
The mixed assignments $(T,W)$ and $(W,T)$ give
sunset amplitude
$\delta F_2^{\text{sun}} = 3/c + 9/(2c) = 15/(2c)$
from the quartic vertex values and propagator products
$\eta^{TT}\eta^{WW} = 6/c^2$, divided by
$|\Aut| = 8$.

\emph{Theta \textup{(}graph F\textup{)}.}
Three edges between two genus-$0$ trivalent vertices.
The mixed-channel assignments involve the coupling
$C_{TWW} = c$ (the gravitational structure constant from
the $WWT$ OPE); odd-parity couplings $C_{TTW} = C_{WWW} = 0$
by the $\mathbb{Z}_2$ parity $W \mapsto -W$.
The theta amplitude with mixed channels gives
$\delta F_2^{\theta} = 1/16$.

\emph{Bridge-loop \textup{(}graph E\textup{)}.}
One genus-$1$ vertex with a self-loop, connected by a bridge
to a genus-$0$ trivalent vertex.
The genus-$1$ vertex contributes $\kappa_j/24$, and the
trivalent coupling $C_{iij}$ with propagators
$\eta^{ii}\eta^{jj}$ gives
$\delta F_2^{\text{bl}} = 21/(4c)$.

Summing:
\begin{equation}\label{eq:w3-genus2-cross-result}
 \delta F_2^{\cross}(\Walg_3)
 \;=\;
 \frac{15}{2c}
 + \frac{1}{16}
 + \frac{21}{4c}
 \;=\;
 \frac{c + 204}{16c}
 \qquad\text{\textup{(}ALL-WEIGHT\textup{)}}.
\end{equation}
The full genus-$2$ free energy is therefore
\begin{equation}\label{eq:w3-genus2-full}
 F_2(\Walg_3)
 \;=\;
 \underbrace{\frac{7c}{6912}}_{\kappa \cdot \lambda_2^{\FP}
 \;\text{\textup{(}UNIFORM-WEIGHT\textup{)}}}
 \;+\;
 \underbrace{\frac{c + 204}{16c}}_{\delta F_2^{\cross}
 \;\text{\textup{(}ALL-WEIGHT\textup{)}}}\,.
\end{equation}
\end{computation}

\begin{proposition}[Explicit $\delta F_2^{\cross}(\Walg_3)$; wave-14
anchor]\label{prop:delta-f-cross-w3-g2}
For the $\Walg_3$ algebra at genus~$2$,
\[
 \delta F_2^{\cross}(\Walg_3)
 \;=\;
 \frac{c + 204}{16c}
 \qquad\text{\textup{(}ALL-WEIGHT\textup{)}},
\]
computed as a graph sum over the seven stable graphs of
$\overline{\cM}_{2,0}$ with rank-$2$ channel assignments.
Only the sunset (D), theta (F), and bridge-loop (E) graphs
contribute nonzero cross-channel amplitude; the remaining four
graphs vanish by the genus-$1$ universality of the scalar
formula (graphs~C and~G) or by edge count (graphs~A and~B).
This proposition anchors the wave-14 scope of Theorem~C(2) in
Vol~I: the multi-weight correction is explicit at
$(\Walg_3, g=2)$ and is the first instance where
$\kappa \cdot \lambda_g^{\mathrm{FP}}$ fails as the complete
genus-$g$ free energy for a multi-generator chirally Koszul
algebra.
\end{proposition}

\begin{proof}
The sunset, theta, and bridge-loop amplitudes are computed in
Computation~\ref{comp:w3-genus2}.  Summing:
\[
 \frac{15}{2c} + \frac{1}{16} + \frac{21}{4c}
 \;=\;
 \frac{120}{16c} + \frac{c}{16c} + \frac{84}{16c}
 \;=\;
 \frac{c + 204}{16c}. \qedhere
\]
\end{proof}

\begin{remark}[Why the scalar formula fails]\label{rem:why-scalar-fails}
The cross-channel correction $\delta F_2^{\cross} = (c{+}204)/(16c)$
is nonzero for all $c > 0$.
At large $c$, $\delta F_2^{\cross} \to 1/16$, while the scalar
part $\kappa \cdot \lambda_2^{\FP} = 7c/6912 \to \infty$.
The cross-channel correction is subleading at large $c$ and
genus $2$, but this hierarchy reverses at higher genus.
At $c = 204$, the cross-channel correction
$\delta F_2 = 408/(16 \cdot 204) = 1/8$ is comparable to the
scalar part $7 \cdot 204/6912 \approx 0.207$.
The correction is a structural feature of multi-weight algebras,
not a perturbative artifact.
\end{remark}


% ================================================================
% Section 4: W_3 at genus 3
% ================================================================
\section{$\Walg_3$ at genus $3$: forty-two stable graphs}
\label{sec:w3-genus3}

\subsection{The genus-$3$ stable graph census}

The moduli space $\overline{\cM}_{3,0}$ has $42$ stable
graph types, stratified by vertex count:
$1$ one-vertex graph, $9$ two-vertex, $14$ three-vertex,
$11$ four-vertex, and $7$ five-vertex.
For the rank-$2$ algebra $\Walg_3$, each graph $\Gamma$
admits $2^{|E(\Gamma)|}$ channel assignments.

Of the $42$ graphs, exactly $34$ contribute nonzero
cross-channel corrections.
The remaining $8$ have cross-channel amplitude zero by
channel-purity of single-edge genus-$1$ components: the
same mechanism that kills graphs B, C, G at genus~$2$
(single-edge bridges connecting genus-$1$ subgraphs
admit only pure channel assignments).

\subsection{The genus-$3$ computation}

\begin{computation}[$\Walg_3$ genus-$3$ cross-channel]\label{comp:w3-genus3}
The genus-$3$ cross-channel correction for $\Walg_3$ is
computed by summing over all $42$ stable graphs
of $\overline{\cM}_{3,0}$, with all $2^{|E(\Gamma)|}$
channel assignments per graph.
The result:
\begin{equation}\label{eq:w3-genus3-result}
\boxed{
 \delta F_3^{\cross}(\Walg_3)
 \;=\;
 \frac{5c^3 + 3792\, c^2
 + 1{,}149{,}120\, c
 + 217{,}071{,}360}{138{,}240\, c^2}
}
\qquad\text{\textup{(}ALL-WEIGHT\textup{)}}.
\end{equation}
All four numerator coefficients are positive, so
$\delta F_3^{\cross}(\Walg_3) > 0$ for all $c > 0$.
\end{computation}

\begin{proof}
Enumerating the $42$ stable graphs of
$\overline{\cM}_{3,0}$ with all channel assignments gives the
full genus-$3$ graph sum.
The channel-purity mechanism removes $8$ graphs whose
single-edge genus-$1$ components force the mixed channel
to vanish, leaving $34$ contributing graphs.
Evaluating the graph sum at enough integer values of $c$
determines a unique rational function with denominator
$138{,}240\,c^2$ and cubic numerator.
Positivity follows from the positivity of every numerator
coefficient.
\end{proof}

\subsection{Cross-channel dominance at genus~$3$}

The large-$c$ asymptotic of $\delta F_3^{\cross}$ is
\begin{equation}\label{eq:w3-genus3-large-c}
 \delta F_3^{\cross}(\Walg_3)
 \;\sim\;
 \frac{c}{27{,}648}
 \qquad (c \to \infty),
\end{equation}
which is \emph{linear} in $c$, in contrast to the genus-$2$
correction which approaches the constant $1/16$.
The net $c$-degree of $\delta F_g^{\cross}$ increases with
genus: degree $0$ at $g = 2$, degree $1$ at $g = 3$.

The ratio to the scalar part reveals a dramatic shift:
\begin{equation}\label{eq:w3-genus3-ratio}
 \frac{\delta F_3^{\cross}(\Walg_3)}%
 {\kappa(\Walg_3) \cdot \lambda_3^{\FP}}
 \;\xrightarrow{\;c \to \infty\;}
 \frac{42}{31}
 \;\approx\; 1.35.
\end{equation}
At genus $3$ the cross-channel correction is already
$35\%$ larger than the scalar part in the large-$c$ regime.
For comparison, at genus $2$ the ratio
$\delta F_2^{\cross}/(\kappa \cdot \lambda_2^{\FP})
\to 0$ as $c \to \infty$: the scalar part dominates.
The crossover occurs between genus $2$ and genus $3$.

\begin{proposition}[Cross-channel dominance]\label{prop:cross-dominance}
The full genus-$g$ free energy of $\Walg_3$ satisfies
\textup{(ALL-WEIGHT)}:
\begin{equation}\label{eq:all-weight-full}
 F_g(\Walg_3)
 \;=\;
 \kappa(\Walg_3) \cdot \lambda_g^{\FP}
 + \delta F_g^{\cross}(\Walg_3).
\end{equation}
The ratio
$\delta F_g^{\cross}/(\kappa \cdot \lambda_g^{\FP})$
grows at least factorially with genus, because
$\kappa \cdot \lambda_g^{\FP}
= O(c \cdot (2\pi)^{-2g})$ decays exponentially in $g$
while $\delta F_g^{\cross}$ is bounded below by a
polynomial in $c$ whose degree grows with $g$.
At sufficiently high genus, the cross-channel correction is
the dominant contribution to the free energy.
\end{proposition}

\begin{proof}
By the Faber--Pandharipande formula,
$\lambda_g^{\FP}
= (2^{2g-1}{-}1)|B_{2g}|/(2^{2g-1}(2g)!)
= O(1/(2\pi)^{2g})$ (Bernoulli asymptotics),
so $\kappa \cdot \lambda_g^{\FP}$ decays exponentially
in $g$.
The cross-channel correction sums over $O(C^g)$ graphs (the
number of stable graphs of $\overline{\cM}_{g,0}$ grows
exponentially), with graph amplitudes bounded polynomially
in $c$.
The resulting growth rate exceeds
$\kappa \cdot \lambda_g^{\FP}$ factorially.
\end{proof}

\begin{remark}[The genus tower]\label{rem:genus-tower}
The explicit data through genus $4$:
\begin{center}
\small
\renewcommand{\arraystretch}{1.4}
\begin{tabular}{clccc}
\toprule
$g$ & Formula for
 $\delta F_g^{\cross}(\Walg_3)$ \textup{(ALL-WEIGHT)} &
 Large-$c$ &
 Graphs &
 $\displaystyle\lim_{c \to \infty}
 \frac{\delta F_g^{\cross}}%
 {\kappa \cdot \lambda_g^{\FP}}$ \\
\midrule
$2$ &
 $\dfrac{c + 204}{16\, c}$ &
 $\dfrac{1}{16}$ &
 $3$ of $7$ &
 $0$ \\[10pt]
$3$ &
 $\dfrac{5c^3 {+} 3792\, c^2 {+}
 1149120\, c {+} 217071360}%
 {138240\, c^2}$ &
 $\dfrac{c}{27648}$ &
 $34$ of $42$ &
 $\dfrac{42}{31} \approx 1.35$ \\[10pt]
$4$ &
 $\dfrac{287\, c^4 {+} 268881\, c^3
 {+} \cdots}{17418240\, c^3}$ &
 $\dfrac{41\, c}{2488320}$ &
 $379$ total &
 $\dfrac{9184}{381} \approx 24.1$ \\
\bottomrule
\end{tabular}
\end{center}
The pattern: the large-$c$ limit of the ratio
$\delta F_g^{\cross}/(\kappa \cdot \lambda_g^{\FP})$ is
$0$, $42/31$, $9184/381$ at genera $2$, $3$, $4$.
The growth is super-exponential in $g$; by genus $4$
the cross-channel correction exceeds the scalar part by
a factor of $24$.
\end{remark}


% ================================================================
% Section 5: The instanton hierarchy
% ================================================================
\section{The instanton hierarchy}\label{sec:instanton}

The scalar genus expansion
$F_g^{\scal} = \kappa \cdot \lambda_g^{\FP}$ (UNIFORM-WEIGHT)
has a completely explicit resurgent structure, determined by
the closed form
\[
 \sum_{g \geq 1} F_g^{\scal}\, \hbar^{2g}
 \;=\;
 \kappa\!\left(
 \frac{\hbar/2}{\sin(\hbar/2)} - 1
 \right).
\]
The poles at $\hbar = 2\pi n$ for $n \geq 1$ produce Borel
singularities at $\xi_n = (2\pi n)^2$, and the leading
instanton action is $A_{\scal} = (2\pi)^2$.

\begin{proposition}[Scalar Borel transform]\label{prop:scalar-borel}
The Borel transform of the scalar genus series in the
variable $u = \hbar^2$,
\[
 \mathcal{B}[Z](\xi)
 \;:=\;
 \sum_{g \geq 1}\,
 \frac{F_g^{\scal}}{(g-1)!}\,\xi^{g-1},
\]
is an entire function of $\xi$
\textup{(}UNIFORM-WEIGHT\textup{)}.
The Borel coefficients satisfy
\[
 \frac{|F_g^{\scal}|}{(g-1)!}
 \;\sim\;
 \frac{2\,|\kappa|}{(2\pi)^{2g}\,(g-1)!},
\]
which decays superexponentially in $g$.
The radius of convergence is infinite.
\end{proposition}

\begin{proof}
The Bernoulli asymptotics $|B_{2g}| \sim 2(2g)!/(2\pi)^{2g}$
give $|F_g^{\scal}| = O(1/(2\pi)^{2g})$.
Dividing by $(g-1)!$ produces superexponential decay.
The root test gives infinite radius.
\end{proof}

\begin{proposition}[Universal instanton action and trans-series]\label{prop:instanton-hierarchy}
Under the uniform-weight hypothesis, the shadow partition
function has the trans-series form
\textup{(}UNIFORM-WEIGHT\textup{)}:
\begin{equation}\label{eq:transseries}
 Z^{\mathrm{full}}(\hbar)
 \;=\;
 \sum_{n \geq 0}\,
 \sigma^n\,
 e^{-n \cdot A_{\scal}/\hbar^2}\,
 Z^{(n)}(\hbar),
\end{equation}
where $A_{\scal} = (2\pi)^2$ is the universal instanton
action, $\sigma$ is the trans-series parameter, and
$Z^{(0)} = Z^{\mathrm{sh}}_{\scal}$ is the perturbative
sector.
The leading Stokes multiplier is
$\mathfrak{S}_1 = -4\pi^2\,\kappa(\cA)\,i$.

For multi-weight algebras $(\Walg_3, \Walg_4, \ldots)$:
the scalar instanton action $A_{\scal} = (2\pi)^2$ still
governs the scalar free energy
$F_g^{\diag} = \kappa \cdot \lambda_g^{\FP}$
\textup{(UNIFORM-WEIGHT)},
but the cross-channel correction
$\delta F_g^{\cross}$ \textup{(ALL-WEIGHT)}
carries its own large-order asymptotics that are not
controlled by a single universal instanton action.
The polynomial degree of $\delta F_g^{\cross}$ in $c$ grows
with genus, so the cross-channel asymptotics exhibit
$c$-dependent growth.
For $\Walg_3$: the ratio
$\delta F_3^{\cross}/\delta F_2^{\cross} \sim c/1728$
at large $c$, growing without bound.
The cross-channel genus expansion grows
\emph{faster} than the scalar expansion at large $c$.
\end{proposition}

\begin{proof}
The trans-series structure follows from the closed-form
expression for the scalar partition function, whose poles at
$\hbar = 2\pi n$ produce the Borel singularities.
The instanton action $A_{\scal} = \xi_1 = (2\pi)^2$ depends
only on the $\hat{A}$-genus normalization, not on the algebra.
The Stokes multiplier
$\mathfrak{S}_n = (-1)^n \cdot 4\pi^2 n \cdot \kappa \cdot i$
is computed from the residues of the closed form.

For the cross-channel part, the explicit data gives:
at genus $2$, $\delta F_2^{\cross}(\Walg_3) = O(1)$
at large $c$; at genus $3$,
$\delta F_3^{\cross}(\Walg_3) = O(c)$.
The ratio $\delta F_3^{\cross}/\delta F_2^{\cross}$
is a rational function of $c$ with positive leading term:
\[
 \frac{\delta F_3^{\cross}}{\delta F_2^{\cross}}
 \;\sim\;
 \frac{c/27648}{1/16}
 \;=\;
 \frac{c}{1728}
 \qquad (c \to \infty).
\]
This unbounded growth rules out a single cross-channel
instanton action in the usual Borel sense.
\end{proof}

\begin{remark}[Gevrey classification]\label{rem:gevrey}
The scalar genus expansion is Gevrey-$1$: the ratio
$F_{g+1}^{\scal}/F_g^{\scal} \sim (2g+1)(2g+2)/(2\pi)^2$
grows quadratically in $g$, confirming factorial divergence with
instanton action $A_{\scal} = (2\pi)^2$ (UNIFORM-WEIGHT).
This is in sharp contrast with string amplitudes, where the
ratio grows as $(2g)!$ and the Borel transform has finite
radius.
The shadow expansion has an \emph{entire} Borel transform
because the Bernoulli decay $1/(2\pi)^{2g}$ kills the
factorial growth.

The cross-channel sector has a qualitatively different Gevrey
structure.
The polynomial-degree growth of $\delta F_g^{\cross}$ in $c$
means that at fixed $c$ the cross-channel series is
factorially divergent (Gevrey-$1$), but at large $c$ the
effective instanton action
$A_{\cross}(c)$ grows with $c$ (ALL-WEIGHT).
The hierarchy $A_{\cross}(c) > A_{\scal} = (2\pi)^2$ at
large $c$ is the quantitative statement that cross-channel
saddles are heavier than scalar saddles.
\end{remark}


% ================================================================
% Section 6: Universal N-formula
% ================================================================
\section{Universal $N$-formula for
$\delta F_2(\Walg_N)$}\label{sec:universal-N}

\subsection{The gravitational Frobenius algebra}

\begin{definition}[Gravitational Frobenius algebra of $\Walg_N$]\label{def:grav-frobenius}
The principal $\Walg$-algebra $\Walg_N$ has strong generators
$W^{(2)}, W^{(3)}, \ldots, W^{(N)}$ at conformal weights
$2, 3, \ldots, N$ ($N{-}1$ generators; highest weight $N$).
The \emph{gravitational Frobenius algebra} of $\Walg_N$ is
the rank-$(N{-}1)$ Frobenius algebra with
per-channel inverse Zamolodchikov metric
$\eta^{(j)(j)} = j/c$ for $j = 2, \ldots, N$
(from the leading-pole normalization $\eta_{(j)(j)} = c/j$)
and structure constants
$C^{\grav}_{(i)(j)(k)} = c \cdot \delta_{ijk}^{\mathrm{even}}$,
nonzero only when the number of odd-weight indices is even
(reflecting the $\mathbb{Z}_2$ parity of odd-spin generators).
\end{definition}

\subsection{The closed-form formula}

\begin{proposition}[Universal gravitational cross-channel formula for
$\Walg_N$]\label{prop:universal-N-formula}
\begin{enumerate}[label=\textup{(\roman*)}]
\item \emph{Closed-form formula.}
 The genus-$2$ gravitational cross-channel correction is
 \textup{(ALL-WEIGHT)}:
 \begin{equation}\label{eq:universal-N-formula}
 \boxed{
 \delta F_2^{\grav}(\Walg_N, c)
 \;=\;
 \frac{(N{-}2)(N{+}3)}{96}
 \;+\;
 \frac{(N{-}2)(3N^3 + 14N^2 + 22N + 33)}{24c}.
 }
 \end{equation}
 Equivalently,
 \begin{equation}\label{eq:universal-N-formula-combined}
 \delta F_2^{\grav}(\Walg_N, c)
 \;=\;
 \frac{(N{-}2)\bigl[c(N{+}3)
 + 4(3N^3 + 14N^2 + 22N + 33)\bigr]}{96c}.
 \end{equation}

\item \emph{Vanishing characterization.}
 $\delta F_2^{\grav}(\Walg_N, c) = 0$ for all $c > 0$
 if and only if $N = 2$
 \textup{(}Virasoro: uniform-weight, no cross-channel
 correction\textup{)}.

\item \emph{Strict positivity.}
 For $N \geq 3$ and all $c > 0$,
 $\delta F_2^{\grav}(\Walg_N, c) > 0$.

\item \emph{Large-$c$ and large-$N$ asymptotics.}
 Define the leading coefficient
 $B(N) := (N{-}2)(N{+}3)/96$
 and the subleading coefficient
 $A(N) := (N{-}2)(3N^3 + 14N^2 + 22N + 33)/24$,
 so that
 $\delta F_2^{\grav} = B(N) + A(N)/c$
 \textup{(ALL-WEIGHT)}.
 At large $N$:
 \[
 B(N) \;\sim\; \frac{N^2}{96}, \qquad
 A(N) \;\sim\; \frac{N^4}{8}.
 \]

\item \emph{Explicit specializations.}
 \begin{align}
 \Walg_3\colon\quad
 \delta F_2^{\grav}
 &= \frac{c + 204}{16c}
 \quad\textup{(}\textbf{exact}\textup{)},
 \label{eq:grav-cross-W3-explicit}\\[4pt]
 \Walg_4\colon\quad
 \delta F_2^{\grav}
 &= \frac{7c + 2148}{48c}
 \quad\textup{(}lower bound\textup{)},
 \label{eq:grav-cross-W4-explicit}\\[4pt]
 \Walg_5\colon\quad
 \delta F_2^{\grav}
 &= \frac{c + 434}{4c}
 \quad\textup{(}lower bound\textup{)}.
 \label{eq:grav-cross-W5-explicit}
 \end{align}

\item \emph{Exactness and lower-bound property.}
 For $\Walg_3$, the gravitational formula is \textbf{exact}:
 the $\mathbb{Z}_2$ parity $W \mapsto -W$ kills all
 higher-spin exchange couplings ($C_{TTW} = C_{WWW} = 0$).
 For $N \geq 4$, higher-spin exchange structure constants
 are generically nonzero and positive, so
 the full cross-channel correction satisfies
 \begin{equation}\label{eq:lower-bound}
 \delta F_2(\Walg_N, c)
 \;\geq\;
 \delta F_2^{\grav}(\Walg_N, c)
 \qquad\text{\textup{(}ALL-WEIGHT\textup{)}}
 \end{equation}
 with equality if and only if $N \leq 3$.
\end{enumerate}
\end{proposition}

\begin{proof}
The genus-$2$ cross-channel correction is computed by
summing over the three contributing stable graphs
D (sunset), E (bridge-loop), F (theta) from the catalogue
of \S\ref{sec:w3-genus2},
with all mixed-channel assignments
$\sigma \colon E(\Gamma) \to \{2, \ldots, N\}$.

\emph{Sunset (graph D).}
The quartic vertex value
$V_{0,4}(i,i,j,j) = \sum_k C_{iik}\,\eta^{kk}\,C_{jjk}$
depends on $i$ and $j$ only through parity.
For even-weight channels: $V_{0,4} = 2c$.
The sunset cross-channel contribution involves mixed
assignments with propagator products
$\eta^{(i)(i)}\eta^{(j)(j)} = ij/c^2$, divided by
$|\Aut| = 8$.

\emph{Theta (graph F).}
The trivalent vertex $C_{(i)(j)(k)}$ is nonzero only for
even-parity triples.
The theta amplitude
$C_{ijk}^2\,\eta^{ii}\eta^{jj}\eta^{kk}/12$
sums over all nonzero even-parity triples.

\emph{Bridge-loop (graph E).}
The genus-$1$ vertex contributes $\kappa_j/24$
(genus-$1$ universality), and the genus-$0$ trivalent
vertex contributes $C_{iij}$.
The amplitude is
$C_{iij}\,\eta^{ii}\eta^{jj}\kappa_j/48$.

The closed-form coefficients
$B(N) = (N{-}2)(N{+}3)/96$ and
$A(N) = (N{-}2)(3N^3 + 14N^2 + 22N + 33)/24$
are obtained by Newton interpolation from the graph sum
evaluated at $N = 3, 4, \ldots, 11$ and verified at
$5$ or more values of $c$ per $N$.
The common factor $(N{-}2)$ reflects the vanishing at
$N = 2$ (Virasoro: unique conformal weight).

Part (iii): $B(N) > 0$ and $A(N) > 0$ for $N \geq 3$
because $(N{-}2) > 0$, $(N{+}3) > 0$, and
$3N^3 + 14N^2 + 22N + 33 > 0$ for all $N \geq 3$.

Part (vi): for $\Walg_3$, the only even-parity triples
involving the $W$-channel are $(T,W,W)$ and permutations.
The cross-coupling $C_{TWW}^{\grav} = c$ accounts for all
gravitational exchange; $C_{TTW} = C_{WWW} = 0$ eliminates
all higher-spin exchange.
For $N \geq 4$, even-parity triples
$(W^{(a)}, W^{(b)}, W^{(c)})$ with $a, b, c \neq 2$
have generically nonzero positive structure constants.
\end{proof}

\begin{remark}[$N$-degree pattern]\label{rem:n-degree}
The polynomial degree in $N$ of the coefficient of $c^{-j}$
in $\delta F_g^{\grav}(\Walg_N)$ satisfies
\begin{equation}\label{eq:n-degree-pattern}
 \deg_N\!\bigl(e_{g,-j}(N)\bigr)
 \;=\; 2j + g
 \qquad (j = 0, 1, \ldots, g{-}1).
\end{equation}
At genus $2$:
$e_{2,-1}(N) = A(N)/D_2$ has $\deg_N A = 4 = 2{\cdot}1+2$;
$e_{2,0}(N) = B(N)/D_2$ has $\deg_N B = 2 = 2{\cdot}0+2$.
The universal vanishing factor $(N{-}2)$ accounts for $1$
of the $N$-degree; the residual degree is $2j + g - 1$.
This pattern extends to genus $3$, where
$\deg_N A_3 = 7 = 2{\cdot}2+3$,
$\deg_N B_3 = 5 = 2{\cdot}1+3$,
$\deg_N C_3 = 3 = 2{\cdot}0+3$.
\end{remark}

\begin{remark}[Spectral curve obstruction]\label{rem:spectral-curve}
No spectral curve with at most two branch points reproduces
the cross-channel tower
$\{\delta F_g^{\cross}(\Walg_3)\}_{g \geq 2}$.
The Newton identity test: define Bernoulli-normalized power sums
$\sigma_g := \delta F_g/(B_{2g}/(2g(2g{-}2)))$.
If $\delta F_g$ arose from topological recursion on a spectral
curve with $\leq 2$ branch points, the $\sigma_g$ would satisfy
a linear recurrence of order $\leq 2$.
The recurrence holds at $g = 2, 3$ but fails at $g = 4$ with
relative error $18$--$27\times$.
The correct framework is CohFT-weighted topological recursion
on the $A_2$ Frobenius manifold.
\end{remark}


% ================================================================
% Section 7: The E_1 interpretation
% ================================================================
\section{The $E_1$ interpretation: cross-channel
dominance as $E_1$ primacy}\label{sec:e1-interpretation}

\subsection{The ordered bar resolves all channels}

The primitive object of the modular Koszul programme is the
ordered bar complex
$B^{\ord}(\cA) = T^c(s^{-1}\bar{\cA})$
with deconcatenation coproduct and differential from the
$E_1$-chiral product.
The $E_1$ convolution algebra $\fg^{E_1}_{\cA}$ carries the
classical $r$-matrix $r(z)$ as the genus-$0$ binary collision
residue of the $E_1$-framed MC element $\Theta^{E_1}_{\cA}$.
The symmetric bar $B^{\Sigma}(\cA)$ is the
$\Sigma_n$-coinvariant shadow, and the modular convolution
algebra $\fg^{\mathrm{mod}}_{\cA}$ is the image of the lossy
coinvariant projection
$\av \colon \fg^{E_1}_{\cA} \to \fg^{\mathrm{mod}}_{\cA}$.

The five theorems of the programme extract the
$\Sigma_n$-invariant content of the ordered bar.
The modular characteristic is
$\kappa(\cA) = \av(r(z))$ at degree $2$ in the abelian and scalar
families; for non-abelian affine Kac--Moody,
$\kappa(V_k(\fg)) = \av(r(z)) + \dim(\fg)/2$.
The scalar formula
$F_g = \kappa \cdot \lambda_g^{\FP}$ (UNIFORM-WEIGHT) is the
statement that the coinvariant projection loses no genus
information for uniform-weight algebras.

\subsection{The cross-channel correction as averaging loss}

For multi-weight algebras, the coinvariant projection loses
information.
The ordered bar $B^{\ord}(\cA)$ resolves all $r$ weight
channels simultaneously: each tensor factor in
$T^c(s^{-1}\bar{\cA})$ carries a channel label, and the
deconcatenation coproduct preserves channel assignments.
The stable-graph expansion of the MC element $\Theta^{E_1}_{\cA}$
computes the full genus-$g$ free energy with all channel
data intact: every edge of every stable graph carries a
definite channel label, and the graph sum over
$\overline{\cM}_{g,0}$ includes all $r^{|E(\Gamma)|}$
channel assignments.

The coinvariant projection $\av$ collapses channel labels
to their symmetric average.
In the $E_1$ picture, the diagonal channel assignments
(all edges in the same channel) are $\Sigma_n$-invariant
and survive the projection; the mixed channel assignments
break the permutation symmetry and are killed.
The cross-channel correction $\delta F_g^{\cross}$ is
precisely the content of $\ker(\av)$ at genus $g$: it is
the information about the ordered bar complex that the
coinvariant projection discards.

\begin{proposition}[$E_1$ characterization of the cross-channel correction]\label{prop:e1-characterization}
The cross-channel correction decomposes as
\textup{(ALL-WEIGHT)}:
\begin{equation}\label{eq:e1-decomposition}
 \delta F_g^{\cross}(\cA)
 \;=\;
 F_g^{E_1}(\cA) - F_g^{\Sigma}(\cA),
\end{equation}
where $F_g^{E_1}(\cA)$ is the genus-$g$ free energy computed
from the ordered bar complex
\textup{(}with all channel assignments\textup{)}, and
$F_g^{\Sigma}(\cA) = \kappa(\cA) \cdot \lambda_g^{\FP}$
is the coinvariant projection
\textup{(}UNIFORM-WEIGHT scalar lane\textup{)}.
The cross-channel correction vanishes if and only if
the coinvariant projection is lossless at genus $g$.
\end{proposition}

\subsection{Quantitative vindication of $E_1$ primacy}

The data of \S\ref{sec:w3-genus3} gives a quantitative
measure of $E_1$ primacy.
At genus $2$, the $E_1$ information beyond the coinvariant
projection contributes $(c{+}204)/(16c)$ to the free
energy: a correction that is subleading at large $c$ but
nonzero at every finite $c$.
At genus $3$, the ordered-bar information already dominates:
the ratio $\delta F_3^{\cross}/F_3^{\Sigma} \to 42/31$
at large $c$ means that more than half the genus-$3$ free
energy lives in $\ker(\av)$, invisible to the symmetric bar.
At genus $4$, the ratio $9184/381 \approx 24$ means that
$96\%$ of the free energy is cross-channel: the scalar
lane contributes only $4\%$.

This is the precise quantitative content of $E_1$ primacy
at multi-weight.
The ordered bar $B^{\ord}(\cA)$ is not merely the technically
convenient resolution of $B^{\Sigma}(\cA)$; it is the primary
object that carries most of the genus-$g$ information.
The symmetric bar, and the scalar formula that comes with it,
captures only the tip of the genus expansion.
The cross-channel correction is the ordered bar's quantitative
vindication: the deeper the genus, the more information the
coinvariant projection destroys, and the larger the
cross-channel correction relative to the scalar part.

\begin{remark}[Irreducible bivariance]\label{rem:bivariance}
The generating function
$\Delta(c, \hbar)
:= \sum_{g=2}^{\infty}
\delta F_g^{\cross}(\Walg_N, c)\,\hbar^{2g}$
(ALL-WEIGHT)
admits no closed-form expression in one variable.
For $\Walg_3$, it is an \emph{irreducibly bivariate}
function of $(c, \hbar)$: it does not factor as
$f(\hbar) \cdot h(c)$, it is not generated by an
$\hat{A}$-genus ansatz, and the numerator polynomial
$P_g(c)$ in
$\delta F_g = P_g(c)/(D_g\,c^{g-1})$
changes qualitative structure with each genus.
Five independent obstructions establish this:
inhomogeneous $c$-scaling ($c$-power window widens with
genus), super-linear ratio growth (consecutive ratios are
non-polynomial in $c$), non-separability (cross-ratio test
fails), $\hat{A}$-genus obstruction
($\delta F_3/(\delta F_2)^2$ varies with $c$), and
irreducible numerators ($P_2$, $P_3$, $P_4$ admit no uniform
Euler-product factorization).
\end{remark}

\begin{remark}[Free fields as unique survivors]\label{rem:free-field-unique}
The free-field algebras (Heisenberg, $\beta\gamma$, $bc$,
lattice) are the unique multi-weight algebras for which the
coinvariant projection is lossless:
$\delta F_g^{\cross} = 0$ for all $g \geq 1$.
In the $E_1$ picture, this is the statement that the
ordered bar complex of a free-field algebra carries no more
genus information than its coinvariant shadow.
The triple redundancy of the proof
(shadow-tower collapse, off-diagonal metric elimination,
ghost-number conservation) expresses three independent
projections of the same geometric fact: the quadratic
free-field OPE factorizes the propagator through weight
sectors, and the ordered bar complex accordingly
decomposes as a product of single-channel bars.
Products have no cross-terms.
\end{remark}


% ================================================================
% References
% ================================================================
\begin{thebibliography}{99}

\bibitem{BD}
A.~Beilinson and V.~Drinfeld,
\emph{Chiral algebras},
AMS Colloquium Publications \textbf{51}, 2004.

\bibitem{CG}
K.~Costello and O.~Gwilliam,
\emph{Factorization algebras in quantum field theory},
Vols.~1--2, Cambridge University Press, 2017, 2021.

\bibitem{DNP25}
T.~Dimofte, N.~Nekrasov, and S.~Paquette,
Line operators from holomorphic-topological field theory,
\emph{Preprint}, 2025.

\bibitem{Ecalle81}
J.~\'Ecalle,
\emph{Les fonctions r\'esurgentes}, Vols.~1--3,
Publ.\ Math.\ d'Orsay, 1981.

\bibitem{FP}
C.~Faber and R.~Pandharipande,
Hodge integrals and moduli of curves,
\emph{Inv.\ Math.}\ \textbf{163} (2006), 25--87.

\bibitem{FLM88}
I.~Frenkel, J.~Lepowsky, and A.~Meurman,
\emph{Vertex operator algebras and the Monster},
Academic Press, 1988.

\bibitem{FB04}
E.~Frenkel and D.~Ben-Zvi,
\emph{Vertex algebras and algebraic curves},
2nd ed., AMS, 2004.

\bibitem{FR94}
I.~Frenkel and N.~Reshetikhin,
Quantum affine algebras and holonomic difference equations,
\emph{Comm.\ Math.\ Phys.}\ \textbf{146} (1994), 1--60.

\bibitem{KZ84}
V.~Knizhnik and A.~Zamolodchikov,
Current algebra and Wess--Zumino model in two dimensions,
\emph{Nuclear Phys.\ B} \textbf{247} (1984), 83--103.

\bibitem{Lor26}
R.~Lorgat,
\emph{Modular Koszul duality},
In preparation, 2026.

\bibitem{LV12}
J.-L.~Loday and B.~Vallette,
\emph{Algebraic operads},
Grundlehren der mathematischen Wissenschaften \textbf{346},
Springer, 2012.

\bibitem{PR18}
T.~Proch\'azka and M.~Rap\v{c}\'ak,
Webs of W-algebras,
\emph{JHEP} \textbf{2018}:11 (2018), 109.

\bibitem{Zam85}
A.~Zamolodchikov,
Infinite additional symmetries in two-dimensional conformal
field theory,
\emph{Theor.\ Math.\ Phys.}\ \textbf{65} (1985), 1205--1213.

\end{thebibliography}

\end{document}
